%=============================================================================
\chapter{Introduction}\label{chap:introduction}
%=============================================================================

\section{Definition and Characteristics of Swarm Robotics}\label{sec:definition}

Swarm robotics is an interdisciplinary approach inspired by social insects (such
as ants, bees, and termites) and by groups of vertebrates (such as bird flocks
and fish schools), in which a large number of simple robots achieve complex group
behaviors through local interactions. Swarm robotics is concerned with
engineering large-scale collective behaviors that, by exploiting the
self-organization of simple agents, yield robust, scalable, and flexible systems
\cite{dias2021swarm,tan2013research}. It has become the most common way of
designing collective behaviors for groups of robots \cite{brambilla2013swarm}.
As shown in \cref{fig:swarm-position}, swarm robotics is a subset of multi-robot
systems.

\begin{figure}[htbp]
  \centering
  \tikzfigure[0.55\linewidth]{swarm-position}
  \caption{The position of swarm robotics within mobile and multi-robot systems.}
  \label{fig:swarm-position}
\end{figure}

\subsection{Definition of Swarm Robotics}

Swarm robotics is a way of coordinating multi-robot systems that consist of large
populations of relatively simple, physically embodied robots. Its central design
problem is to choose the behavior of the individual robots so that a desired
macroscopic behavior emerges from local robot--robot and robot--environment
interactions, without centralized control \cite{brambilla2013swarm,sahin2005swarm}.
The design approach focuses on scalability, flexibility, and robustness, and it
relies on local sensing and communication combined with distributed autonomy
\cite{dias2021swarm,tan2013research}.

\subsubsection{Study of Collective Intelligence in Robot Groups}

Swarm robotics studies how a group of simple robots, relying on emergent
collective intelligence, can cooperatively accomplish tasks that are beyond the
capabilities of any single robot \cite{tan2013research,sahin2005swarm}:
\begin{itemize}
  \item \lead{Emergent behavior} Complex behaviors such as aggregation,
        dispersion, and collective flight arise spontaneously from simple local
        rules and give the swarm capabilities that exceed those of the individual.
  \item \lead{Simple agents} Swarm robots are typically homogeneous and of low
        complexity, with minimal on-board storage and computation, and they rely
        on coordination to achieve complex goals collectively.
  \item \lead{Minimal control} The swarm operates with little or no human
        guidance and achieves decentralized autonomy by following simple
        cooperative rules; a group that needs continuous human operation is not
        usually regarded as a swarm \cite{shahzad2023review}.
\end{itemize}

\subsubsection{Merging Swarm Intelligence with Robotics}

Swarm robotics grew out of swarm intelligence and can be read as swarm
intelligence applied to embodied robots: it implements the principles of
decentralized, self-organized problem solving in robotic teams under a fully
distributed control paradigm \cite{dias2021swarm,nguyen2024swarm,hamann2008space}:
\begin{itemize}
  \item \lead{The swarm intelligence paradigm} It employs decentralized,
        self-organized systems to solve complex problems and provides the
        computational techniques and inspiration needed for robust collective
        behavior in robotics.
  \item \lead{Swarm intelligence in applied robotics} Models of coordination
        observed in biological systems are turned into coordination and
        task-allocation algorithms that can be implemented on robot swarms.
  \item \lead{Decentralized coordination} Purely local sensing and
        communication, together with distributed decisions, eliminate single
        points of failure and allow the system to grow to large teams while
        maintaining performance.
\end{itemize}

\subsection{Inspiration from Nature}

Drawing on nature, and based on observations and mathematical models of
biological swarms, this field engineers corresponding coordination and
problem-solving mechanisms for robotics. The success of natural systems in
solving complex tasks in a decentralized manner, from ant colonies and bee
swarms to bird flocks, fish schools and herds of ungulates, provides templates
for robust and adaptive robotic collectives \cite{tan2013research}.

\subsubsection{Biological Swarms as Models}

Biological collectives, from micro-organisms to large vertebrates, show how many
simple agents can reach sophisticated collective outcomes, and the behaviors they
produce are typically robust, flexible, and scalable \cite{dias2021swarm,bonabeau1999swarm}.
\begin{itemize}
  \item \lead{Social insects} Ants, bees, wasps, and termites are the main sources
        of inspiration. Termites build large and intricately ventilated mounds,
        army ants organize foraging raids involving up to 200,000
        workers, and groups of ants carry prey far heavier than any single
        worker could move \cite{bonabeau1999swarm,garnier2007biological}.
  \item \lead{Tasks transferred to robotics} The behaviors that have been carried
        over to robots include foraging (collecting scattered items), nest
        construction, division of labor and task allocation, and path finding;
        desert ants, for example, solve navigation problems by path integration
        \cite{hamann2018swarm,tan2013research}.
  \item \lead{Vertebrate groups} Bird flocks, fish schools, and insect swarms
        provide models of collective motion. Reynolds' \emph{boids} showed that
        three local rules (separation, alignment, and cohesion) are enough to
        reproduce realistic flocking \cite{reynolds1987flocks}.
\end{itemize}
The common lesson of these systems is that collective intelligence emerges
through self-organization: individuals that are simple and interact only
locally can still reach goals far beyond their own capacities
\cite{garnier2007biological}.

\subsubsection{Translating Natural Algorithms to Robotic Control}

Turning a natural behavior into a robot controller usually follows three steps
\cite{bonabeau1999swarm,seebacher2025collective,chen2025survey}:
\begin{enumerate}
  \item \lead{Modeling and simulation} The collective behavior is first studied in
        the natural system, then modeled mathematically, and the model is used
        as the basis for artificial variants; swarm intelligence treats such
        natural phenomena as computational processes from which algorithms can
        be derived.
  \item \lead{Emulating collective features} The resulting algorithms keep the
        features that make the natural swarm work, such as decentralized
        decisions, self-organization, and adaptive responses to the environment,
        and apply them to optimization, control, classification, and routing.
  \item \lead{Bio-inspired algorithms} The best-known outcomes are metaheuristics:
        ant colony optimization (ACO) mimics pheromone trail laying in foraging
        ants and is used for routing and coverage, while particle swarm
        optimization (PSO) mimics the coordinated motion of flocks and schools to
        improve candidate solutions iteratively. Both are reviewed in
        \cref{chap:algorithms}.
\end{enumerate}
Compared with planners that rely on an explicit model of the environment, these
methods let robots navigate complex environments adaptively and support
decentralized decision-making, which is valuable for robustness and distributed
problem solving \cite{chen2025survey}.

\subsection{Key Characteristics of Swarm Robotic Systems}

Swarm robotic systems prioritize autonomy, decentralization, scalability,
robustness, flexibility, and limited individual capabilities, so that high
performance arises from a collective of simple components rather than from a few
complex units \cite{dias2021swarm,brambilla2013swarm}.

\subsubsection{Self-Organization and Decentralization}

Swarm robotics dispenses with centralized control: each robot acts autonomously
based on local sensing and interaction with its neighbors, and global behaviors
emerge from repeated local exchanges \cite{dias2021swarm}.
\begin{itemize}
  \item \lead{Decentralized control} There is no leader or central authority, so
        the vulnerability of a single point of failure disappears.
  \item \lead{Self-organization} Complex global patterns form through local
        interactions among simple robots and adapt to environmental cues and
        task demands. Since the macroscopic behavior cannot be programmed
        directly, the swarm engineer designs the individual behaviors so that the
        desired collective behavior emerges from their interactions
        \cite{brambilla2013swarm}.
\end{itemize}

\subsubsection{Robustness, Scalability, and Flexibility}

These traits are the main advantages of adopting a swarm approach in dynamic and
uncertain environments that call for redundancy and adaptability
\cite{lin2022review,tan2013research}.
\begin{itemize}
  \item \lead{Robustness} Redundancy and distribution allow operation to continue
        despite disturbances or failures of individuals, so that performance
        degrades gracefully: the loss of some members is compensated for by the
        others.
  \item \lead{Scalability} Enforcing local coordination enables effective
        operation from small groups up to thousands of agents without redesigning
        the software or hardware. The essential condition is that all
        interactions remain local.
  \item \lead{Flexibility} Self-organized structures support rapid re-tasking
        for diverse goals and environments with only minor software changes.
\end{itemize}

\subsubsection{Limited Individual Capabilities (Sensing and Communication)}

Individual robots are simple, small, and low-cost, which makes it possible to
deploy large populations that achieve collective productivity through
cooperation; a single robot is often incapable or inefficient at the main task
on its own \cite{tan2013research}. Local sensing and short-range communication
restrict interactions to neighbors, which is vital for the scalability and
stability of distributed coordination in swarms \cite{tan2013research,dias2021swarm}.

%-----------------------------------------------------------------------------
\section{Distinguishing Swarm Robotics from Other Paradigms}\label{sec:distinction}

Swarm robotics differs from multi-robot systems, sensor networks, and multi-agent
systems in terms of scale, control, and applications. \Cref{tab:comparison}
summarizes these differences \cite{tan2013research}.

\begin{table}[htbp]
  \centering
  \caption{Comparison of swarm robotics with other paradigms \cite{tan2013research}.}
  \label{tab:comparison}
  \small
  \rowcolors{2}{lightbg}{white}
  \begin{tabularx}{\linewidth}{@{}>{\bfseries}p{0.17\linewidth}LLLL@{}}
    \toprule
    \rowcolor{white}
    Feature & \textbf{Swarm robotics} & \textbf{Multi-robot system}
            & \textbf{Sensor network} & \textbf{Multi-agent system} \\
    \midrule
    Population size & Highly variable & Small & Fixed & Small \\
    Control & Decentralized & Centralized / remote & Centralized / remote
            & Centralized / hierarchical \\
    Homogeneity & Homogeneous & Heterogeneous & Homogeneous
                & Homogeneous / heterogeneous \\
    Flexibility & High & Low & Low & Medium \\
    Scalability & High & Low & Medium & Medium \\
    Environment & Unknown & Known / unknown & Known & Known \\
    Motion & Yes & Yes & No & Rare \\
    Applications & Disaster relief, military, hazardous tasks
                 & Transportation, sensing, robot soccer
                 & Surveillance, protection
                 & Resource management, control \\
    \bottomrule
  \end{tabularx}
\end{table}

%-----------------------------------------------------------------------------
\section{Importance of Path Planning in Swarm Robotics}\label{sec:importance}

Path planning is a key function in the navigation of autonomous systems, and it
becomes doubly important in swarms because of the large number of agents,
decentralized control, and strict safety and efficiency requirements. In a swarm
robotic system, the planner is the high-level decision layer that selects the
sequence of waypoints each robot follows, and its coordination logic must produce
safe, reliable, and efficient collective motion under local sensing and
communication and under real-time constraints
\cite{nguyen2024swarm,lin2022review,abujabal2024path}. A general classification of
path planning methods, following \cite{abujabal2024path}, is shown in
\cref{fig:taxonomy}.

\begin{figure}[htbp]
  \centering
  \tikzfigure{taxonomy}
  \caption{Classification of path planning methods (redrawn after \cite{abujabal2024path}).}
  \label{fig:taxonomy}
\end{figure}

\subsection{The Fundamental Challenge in Multi-Robot Systems}

Multi-robot path planning requires integrating collision avoidance, optimality,
dynamic feasibility, and team-level coordination, and it is generally treated as
a combinatorial optimization problem, NP-hard in general, that scales poorly with
team size and environmental complexity \cite{lin2022review}. In the swarm setting,
as the number of agents grows, the interaction graph and the coordination
decisions expand sharply, and the computational and real-time challenges
intensify well beyond those of single-robot planning
\cite{nguyen2024swarm,lin2022review,abujabal2024path}. In swarm robotics, each agent is
responsible for planning its own path. In \cref{fig:architectures}, the
processing units are shown in red; from left to right, the figure depicts
centralized, decentralized, and distributed control, respectively. The goal of
swarm robotics is to achieve the distributed architecture (the last case).

\begin{figure}[htbp]
  \centering
  \tikzfigure[0.82\linewidth]{control-architectures}
  \caption{Control architectures in multi-agent robotics: (a) centralized,
           (b) decentralized, and (c) distributed. Processing units are shown in
           red.}
  \label{fig:architectures}
\end{figure}

\subsubsection{Ensuring Safe, Reliable, and Smooth Motion}

\begin{itemize}
  \item \lead{Safety and reliability} Planning must guarantee safety with respect
        to static and moving obstacles in cluttered spaces while preserving group
        motion; this is critical in search-and-rescue and air-traffic scenarios,
        where failure is unacceptable. Relying solely on local formation or
        flocking breaks down in crowded areas and rough terrain, so explicit
        planning constructs such as safe corridors or tubes are needed to
        maintain cohesion and resolve conflicts
        \cite{lin2022review,abujabal2024path,mao2023optimal,mao2024tube}.
  \item \lead{Smoothness} Because of kinematic and nonholonomic constraints, the
        path must be a smooth curve, and the trajectory that follows it must
        respect bounds on velocity and acceleration in order to be trackable.
        This motivates optimization-based and hierarchical ``plan-then-track''
        designs, in which a trajectory-planning stage turns the geometric path
        into a feasible motion profile. Virtual corridors and tubes and smooth
        trajectory parameterizations (e.g., B-splines and piecewise polynomials)
        are widely used for smooth, collision-free motion in real time
        \cite{lin2022review,abujabal2024path,mao2023optimal}.
\end{itemize}

\subsubsection{Optimal Collision-Free Path Generation}

\begin{itemize}
  \item \lead{Collision avoidance} The planner must guarantee collision-free paths
        with respect to static obstacles, moving obstacles, and other robots, and
        in swarms it usually also maintains inter-agent separation, connectivity,
        and velocity consistency through constraints or cooperative control
        layers. A common mechanism is to enlarge the obstacles and shrink the
        robot to a point, which is formulated in the configuration space and
        simplifies the safety constraints \cite{lin2022review,abujabal2024path}.
  \item \lead{Optimal path selection} Planning searches for an optimal or
        near-optimal path according to one or more criteria, such as path
        length, time, energy, smoothness, and robustness, which often leads to
        multi-objective formulations that trade off efficiency, safety margins,
        and dynamic feasibility \cite{nguyen2024swarm,lin2022review}.
\end{itemize}
In swarms, optimality also encompasses collective measures such as group
cohesion, throughput in narrow passages, and coordination costs under
distributed execution.

\subsection{Specific Challenges of Swarm Path Planning}

Path planning for swarm robots combines the challenges specific to individual
robots with a large scale, limited computational and energy resources, and the
need for decentralized operation, which makes traditional centralized methods
unsuitable for large teams. Consequently, algorithms must emphasize locality,
distribution, and real-time responsiveness while preserving safety and
efficiency guarantees \cite{nguyen2024swarm,lin2022review}.

\subsubsection{Computational Complexity with Large Numbers of Robots}

Optimizing the trajectory of every robot in a swarm is expensive: the
computational complexity of optimizing one parameterized trajectory ranges from
$O\bigl(k(n_t,\varepsilon)\,n_t^{2}\bigr)$ to
$O\bigl(k(n_t,\varepsilon)\,n_t^{3}\bigr)$, where $n_t$ is the number of
trajectory parameters and $k$ is the number of iterations required to reach
accuracy $\varepsilon$, and this cost is paid once per robot. This motivates
specialized structures that curb such growth, such as the virtual tubes of
\cref{sec:virtual-tube} \cite{mao2023optimal}.
\begin{itemize}
  \item \lead{Combinatorial growth} As the numbers of agents and obstacles
        increase, the decision space, the planning cost, and the risk of
        collision all grow rapidly \cite{zhang2023multi}. In multi-agent
        reinforcement learning, the joint action space grows exponentially with
        the number of robots and their degrees of freedom, which is the familiar
        curse of dimensionality \cite{huttenrauch2019deep,kuckling2023recent}.
  \item \lead{Scalability of the architecture} Centralized architectures suffer
        from heavy computational and communication loads and struggle with the
        dynamic behaviors and data traffic of large teams, whereas distributed
        methods spread the decision-making over the agents, reduce the burden on
        any single controller, and better preserve real-time responsiveness
        \cite{lin2022review,abujabal2024path,azzam2023swarm}.
\end{itemize}

\subsubsection{Energy Efficiency and Fault-Tolerance Requirements}

\begin{itemize}
  \item \lead{Energy efficiency} Swarm deployments in remote or harsh
        environments, where recharging is often impossible, require minimizing
        energy through shorter paths, smooth velocity profiles, efficient
        communication, and lightweight computation, so as to maximize endurance
        and mission success \cite{lin2022review,abujabal2024path,nguyen2024swarm}.
  \item \lead{Fault tolerance and robustness} The redundancy of the swarm must
        preserve performance under individual failures through task
        redistribution, and the swarm must remain functional under high failure
        rates in hazardous environments \cite{nguyen2024swarm,lin2022review}.
  \item \lead{Robust communication} The planning and tracking layers must cope
        with communication failures, recover from network partitions, work with
        low bandwidth, and re-synchronize teams, so that collective objectives
        are maintained despite disturbances \cite{abujabal2024path}.
\end{itemize}

\subsubsection{Handling Dynamic Targets and Environments}

Real-world search-and-rescue and field operations take place in large, harsh,
and highly dynamic scenarios, which require planners that adapt online to
unknown maps, moving obstacles, and changing layouts, beyond the reach of static
or purely local heuristics \cite{lin2022review,abujabal2024path,chen2025survey}.
\begin{itemize}
  \item \lead{Mobile and dynamic targets} Swarms must maintain coverage and
        tracking of targets that appear, move unpredictably, or increase in
        number, which requires continuous re-targeting and cooperative
        reassignment of motion goals \cite{li2019twostage,azzam2023swarm}.
  \item \lead{Real-time adaptation} To avoid costly replanning, algorithms
        increasingly rely on incremental repairs of existing paths, hierarchical
        schemes, safe corridors or tubes, and learning in the style of
        centralized training with decentralized execution (CTDE), which update
        paths incrementally and generate smooth, collision-free actions in real
        time \cite{li2021rrt,mao2023optimal,azzam2023swarm}.
\end{itemize}

%-----------------------------------------------------------------------------
\section{Objectives and Structure of the Report}\label{sec:structure}

\subsection{Objectives}

The objective of this report is to provide a structured and rigorous review of
path planning in swarm robotics, with an emphasis on strategies that are based on
swarm intelligence and enabled by artificial intelligence (AI), which unlock
scalable, decentralized systems for real-world applications. In particular, the
report aims to:
\begin{enumerate}
  \item \lead{Establish the foundations} Review the key concepts of swarm
        intelligence, namely self-organization, decentralization, emergent
        behavior, and local interaction, as they apply to robotic systems.
  \item \lead{Analyze control architectures} Categorize the decentralized,
        hierarchical, behavior-based, and hybrid control methods used in swarm
        robotics and their effect on system performance.
  \item \lead{Explore path planning techniques} Explain the principles,
        strengths, and weaknesses of deterministic, AI-based (including
        bio-inspired metaheuristics and learning), and hybrid planners when
        applied to decentralized coordination.
  \item \lead{Address computational challenges} Examine how decentralized
        strategies keep planning tractable and scalable when the number of
        robots is large.
  \item \lead{Evaluate optimization metrics} Detail the multi-objective metrics of
        swarm operation: short, safe, and smooth paths under energy and
        fault-tolerance constraints.
  \item \lead{Identify future directions} Outline promising research avenues,
        including adaptive learning, human--swarm interaction, and operation in
        highly dynamic environments.
\end{enumerate}

\subsection{Structure}

The rest of the report is organized as follows:
\begin{itemize}
  \item \lead{\Cref{chap:fundamentals} -- Fundamental concepts and methods in
        swarm control} Swarm intelligence and self-organization, stigmergy, the
        control architectures (decentralized, hierarchical, behavior-based, and
        hybrid), communication and sensing, and the fundamentals of path
        planning, including objectives, metrics, and global versus local
        strategies.
  \item \lead{\Cref{chap:algorithms} -- Path planning algorithms for swarms} A
        systematic review of deterministic (graph- and sampling-based and
        potential-field), AI-based (bio-inspired metaheuristics such as ACO and
        PSO, as well as learning), and hybrid methods, and of virtual-tube
        planning for large swarms.
  \item \lead{\Cref{chap:challenges} -- Challenges and applications} Dynamic and
        complex environments, scalability and robustness, task allocation and
        cooperative control, and the main application domains of swarm robotics.
  \item \lead{\Cref{chap:conclusion} -- Summary, conclusions, and future work} The
        effectiveness of the approaches is summarized, the advantages of swarm
        robotics are stated, and directions for future research are outlined.
\end{itemize}
